% Ad Atticum 7.2 (ad-atticum-07-02) --- English translation
% Translated by Claude (Anthropic) from the Latin (Perseus).
% Style guide: STYLE.md.

\ciceroLetterOpener{CICERO ATTICO salutem}

\ciceroSection{1} We reached Brundisium on the seventh day before the Kalends of December, enjoying your sailing luck: such a pretty little spondaic breeze blew us from Epirus --- the gentlest Onchesmite. That \textit{[Greek: ending-in-spondees]} hexameter, if you like, you may pass off as your own to anyone you please.

\ciceroSection{2} Your health throws me into the deepest disquiet; for your letters intimate that you are not at all well. And, knowing your constitution as I do, I suspect that there is something rather more violent at work which is forcing you to give ground and all but breaking you. Yet your man Pamphilus tells me that one of the quartans has gone and the other is coming on more lightly. Terentia, who in fact reached the Brundisium gate at the very hour I made port and who met me in the forum, was telling me that Lucius Pontius had said to her at the Trebulan villa that this one too has gone. If that is so, it is what I most pray for, by Hercules, and I trust that your good sense and self-discipline are what have brought it about.

\ciceroSection{3} I come to your letters; of which I received six hundred at one time, each more delightful than the last --- the ones, that is, in your own hand. For I loved the hand of Alexis only because it came so near in resemblance to yours; the letters themselves I did not love, because they told me you were unwell. Since the subject has come up: I left Tiro sick at Patrae --- a young man, as you know, of the soundest character (and add anything you like to that). I have seen no better. So I miss him keenly; and although he seemed not to be in a serious way, all the same I am anxious, and I have my greatest hope in the diligence of Manius Curius, of whom Tiro wrote to me and of whom many have brought me word. Curius himself sensed how much you wanted me to make a friend of him, and on that score I am exceedingly pleased. By Hercules, the urbanity in the man is the kind one easily warms to --- \textit{[Greek: autochthonous]}, native to him. I am bringing back his will sealed with the signets of the three Ciceros and of the praetorian cohort. He has made you his heir of an as and me of a quarter-as, openly. At Actium, on Corcyra, Alexio entertained me lavishly. Quintus Cicero could not be prevented from seeing the Thyamis.

\ciceroSection{4} I am glad your little daughter is a joy to you and that you approve the \textit{[Greek: natural]} character of \textit{[Greek: the bond toward one's children]}. For if this is not natural, then there can be no natural bond between one human being and another --- and when that is taken away, the fellowship of life is taken away with it. Bene eveniat! says Carneades coarsely, and yet more wisely than our Lucius and Patron, who, since they refer everything to themselves, never suppose that anything is done for another's sake; and when they say that a good man must be good lest he suffer harm --- not because the thing is right by nature --- they fail to understand that they are describing the cunning man, not the good one. But this, I think, belongs to those books of mine which you have given me heart for by your praise.

\ciceroSection{5} I return to the matter. How I was waiting for the letter you had given to Philoxenus! For you had written that it contained an account of Pompey's conversation at Naples. Patron delivered it to me at Brundisium; he had picked it up, I imagine, at Corcyra. Nothing could have been more welcome. For it was about the commonwealth, about the opinion that great man holds of my integrity, about the goodwill he showed in the conversation he had concerning the triumph. But the most welcome thing of all was this: I understood that you had gone to him to take the measure of his disposition toward me. That, I say, is what came to me as most welcome.

\ciceroSection{6} As for the triumph, no desire for it ever took hold of me before those most impudent letters of Bibulus --- which were followed by the most lavish supplication possible. From a man by whom, if the deeds had really been done that he wrote of, I should be glad and would back the honour; but as it is, that the man who never put a foot outside his gate so long as the enemy was on this side of the Euphrates should be honoured, while I, in whose army the hopes of his army rested, should not obtain the same --- this is a disgrace to us. To us, I say, joining you to it. So I shall try everything, and, I trust, succeed. If you were in health, certain things would already be settled for me; but, as I trust, you will be.

\ciceroSection{7} About the little debt of Numerius, I am much obliged to you. I am eager to know what Hortensius has been doing, and what Cato is up to --- who, by the way, behaved shabbily toward me out of plain ill will. He bore witness in my favour as to integrity, justice, clemency, and good faith --- which I was not seeking; what I was asking, he refused. So how Caesar, in the letter in which he congratulates me and promises me everything, gloats over the ungrateful Cato's affront to me! And yet this same Cato voted Bibulus twenty days. Forgive me; I cannot bear it and I will not.

\ciceroSection{8} I want to answer all your letters, but it is not necessary; for I shall soon see you. One thing, though, about Chrysippus --- for about the other I am less astonished, a labourer of a man; though even of him nothing could be more disgraceful. But that Chrysippus, whom I willingly looked on and held in honour for some little learning he had, should slip away from the boy without my knowledge! I pass over the many other things I hear, I pass over the thefts; the flight I cannot pass over, which seemed to me the most criminal thing of all. So I have invoked that old formula of Drusus the praetor, as they tell it, against the man who being free refused the same oath: that I have never declared those people free --- especially as there was no one present from whom they might be properly claimed. You will take this as you see fit; I shall agree with you. To one of your most eloquent letters I have not yet replied --- the one in which you discuss the dangers facing the commonwealth. What was I to write back? I was very much shaken. But that I should be in no great fear, the Parthians see to: they have suddenly left Bibulus half-alive.
